# SNAPSHOT_SOURCE: lindicaphxag-tech/cs@0b26a5b0853e3b136a92c17e614a3a1936673621 path=papers/eeg-default-regularisation/paper_vnext/03_discussion_current.tex
\section{Discussion}

The reconstruction supports a different paper from the one suggested by the
historical draft. The strongest result is not a new shrinkage estimator and
not a universal multiplier. It is a validity result: covariance fidelity and
cross-participant discrimination can define different regularisation regimes,
and the choice between those objectives can alter both held-participant
performance and the scientific conclusion drawn from an EEG representation
comparison.

The operator analysis clarifies why this is possible without claiming that
the abstract distinction between covariance estimation and classification is
new. Along the scalar shrinkage path, the standardized evaluation margin is
invariant to the component of \(w_\lambda'\) that merely rescales the
discriminant. Only the component that rotates the discriminant in the
evaluation participant's covariance geometry can change the margin. The
covariance-predictive derivative, by contrast, contains no class-contrast
vector. Hence covariance stationarity provides no general guarantee of
discriminative stationarity. Near an interior covariance optimum, the
covariance loss can be locally flat to first order while the discriminative
margin changes to first order. This local asymptotic statement is conditional:
it does not apply unchanged when the covariance optimum lies on the boundary
of the admissible path. That distinction is empirically relevant here,
because BNCI2014\_002 selects \(\lambda_C=0\) in 13/14 outer folds,
whereas BNCI2015\_001 selects interior values between 0.015 and 0.05. The
Validity Frontier is therefore the primary empirical object in both cohorts;
the local square-root mechanism is a compatible explanation only where its
interior-optimum assumptions hold. The frontier measures how much decision
variation is admitted by a given fidelity tolerance rather than treating
numerical proximity in covariance loss as evidence of decision equivalence.

The EEGMMIDB S021--S040 block establishes that this distinction can have a
large prospective consequence. Source participant-discrimination selection
improved macro AUC by approximately 0.056 over the stronger source
covariance-NLL criterion and by 0.092 over Ledoit--Wolf, with both predeclared
comparisons surviving Holm correction. The delete-one-source reselection
audit is important because it changes the object that appears stable. The
selected numerical strength moves from 0.4 to 0.95 when source composition is
perturbed, whereas the ordering of source objectives retains a positive target
consequence in all eight omissions. The empirical regularity is therefore not
that \(\lambda=0.8\) is the correct EEG value; it is that a
discrimination-matched source objective can occupy a different region from a
covariance-matched objective and can transfer better under this protocol.

The TRUE PRE-DATA external panel strengthens and simultaneously narrows that
claim. BNCI2014\_002 and BNCI2015\_001 differ in acquisition geometry,
channel count, tangent dimension, trial count and session structure, yet the
two source objectives select different regularisation regions in 26/26 outer
folds. The target-effect direction is also positive in both datasets:
source-macro selection exceeds covariance-NLL by +0.0176 and +0.0433 AUC,
respectively, and exceeds Ledoit--Wolf by +0.0229 and +0.0453. However,
neither external dataset passes its own frozen exact sign-flip plus Holm
family. We therefore regard the source-objective separation as strongly
replicated in the studied motor-imagery cohorts and the target-performance
consequence as directionally replicated, not as two additional individually
significant superiority experiments. The distinction matters because the
participant bootstrap intervals and exact randomisation tests answer related
but non-identical finite-sample questions; choosing whichever is more
favourable after seeing the outcome would defeat the purpose of the pre-data
protocol.

The corrected Validity Frontier adds a second form of external replication.
The observed max-minus-min AUC spans exceed 0.02 over the frozen dense
regularisation path in both TRUE PRE-DATA cohorts. Yet a positive percentile
bootstrap lower endpoint for a data-selected maximum span does not by itself
establish a population-level material path effect: even a flat population AUC
path can produce a positive empirical max--min span from sampling noise across
hundreds of lambda points. We therefore report these frontiers as descriptive
post-score measurements and add an explicitly post-outcome simultaneous-path
sensitivity diagnostic, rather than interpreting the original percentile
intervals as proof of materiality. The raw fidelity budget required to
reach a 0.02 span differs substantially: approximately 0.916 per-dimension
NLL-regret units in BNCI2014\_002 and 0.175 in BNCI2015\_001. This
heterogeneity is scientifically useful. It rejects the idea of one universal
critical fidelity tolerance while supporting a more general statement that
covariance fidelity defines a cohort-dependent decision-validity frontier.
The dense-grid subset audit further shows that the observed spans are not an
artifact of one particular grid resolution.

The prospective representation experiment demonstrates why this validity
question matters beyond classifier tuning. The regularisation objective
significantly changes the filter-bank-versus-wideband contrast, but the
direction is the opposite of the historical hypothesis:
discrimination-oriented regularisation strengthens the wide-band advantage
rather than uncovering a hidden filter-bank gain. The matched-dimension PCA
control remains below WIDE231, while the nuisance-dimension and source-budget
interventions do not support a simple monotonic feature-to-sample explanation.
Thus the paper does not rescue its original representation story. It
prospectively falsifies it while preserving the more general conclusion that
regularisation objective can alter representation-level inference. This is a
stronger measurement-validity result than a post-hoc ranking reversal selected
because it supports the original narrative.

The failed algorithmic corrections are similarly informative. Lower-tail
signed selection did not beat correctly nested source tuning; CFCTS was
prospectively inferior; FCTR collapses toward a covariance-fidelity region
that underperformed on S021--S040; the TRUE PRE-DATA mean--tail--fidelity
selector does not satisfy its frozen two-cohort tail promotion gate; and the
post-S021 discriminability-first lexicographic rule is slightly worse than
pure source-macro selection in both external cohorts. These results eliminate
several plausible but unsupported claims of algorithmic novelty. They also
suggest that the difficult part of the problem is not constructing a more
ornate scalar selector. Participant heterogeneity creates a trade-off between
mean transfer benefit and harmed-participant risk that is not resolved by a
single source-only tie-breaker in the present data.

The S061--S080 calibration-budget experiment provides another useful
negative result. Its unlabeled adaptation subtracts one constant
participant-specific feature translation before a fixed linear score.
Because this only adds a constant to every score from that participant,
within-participant ROC ordering cannot change and ROC AUC is algebraically
invariant. The observed zero U-\(k\)/Z0 AUC differences are therefore not
evidence that unlabeled calibration is practically useless. They expose an
endpoint--operator identifiability mismatch. A future practical BCI extension
must use a source-frozen operating threshold or another threshold-sensitive
utility measure; such an extension is necessarily later than the current
outcome and cannot be presented as part of the original prospective claim.

These findings also define the novelty boundary relative to prior work.
Covariance shrinkage, regularised discriminant analysis, participant-wise
cross-validation, Riemannian tangent features and zero-calibration transfer
are established techniques, and recent EEG work has explicitly studied the
effects of regularisation and validation choices. The contribution here is
therefore not the observation that a hyperparameter influences accuracy. It
is the combined, auditable chain from objective-specific operator geometry to
full-domain source conflict, prospective target consequence, source-pool
sensitivity, true-pre-data external replication, uncertainty-aware validity
frontiers and prospective falsification of a representation conclusion under
one information budget.

Several limits remain. The strongest external evidence is currently confined
to motor imagery, and all headline regularisation results use linear
discriminants on covariance/tangent representations. The two TRUE PRE-DATA
external cohorts are independent acquisition strata but still too few for a
general EEG-wide law or a meaningful dataset-level meta-analysis. AUC measures
ranking and does not by itself establish calibrated operating-point utility.
The harmed-participant tail remains non-trivial even when the average effect
is favourable. Finally, the operator derivation identifies a local mechanism
but has not yet been promoted to a target-label-free predictor of who will
benefit; doing so would require a new independently frozen experiment rather
than reusing the current target outcomes.

For these reasons, the present scope is deliberately narrower than the
historical title implied. The evidence supports a BSPC-style
signal-processing and measurement-validity contribution in the studied
cross-subject covariance/tangent discriminant family. A JNE submission becomes
plausible if the final paper emphasizes the neural-engineering consequence of
representation-dependent validity and preserves the unusually strict evidence
chain. TNSRE would still require a separate direct BCI-utility extension,
preferably threshold-sensitive, cross-session/device or online-like. Stronger
wording cannot substitute for that evidence.
